\section*{الفصل \ref{ch:g9:landforms} --- الأشكال التضاريسية}

\begin{solution}{exo:g9:landforms:1}
(أ) الريح. (ب) ماء يذيب \omterm{def:g3:stones-pebbles-sand:rock}{الصخر} (ويرسّب الكالسيت في \omterm{def:g9:landforms:karst}{الكهف}). (ج) \omterm{def:g2:water-on-earth:ocean}{البحر}
(\omterm{def:g9:landforms:beach}{الانجراف الساحلي}). (د) الجليد. (ه) \omterm{def:g2:water-on-earth:ocean}{البحر}. (و) \omterm{def:g1:landscapes-around-us:water}{نهر}.
\end{solution}

\begin{solution}{exo:g9:landforms:2}
(أ) \omterm{def:g9:landforms:uplift}{الرفع}. (ب) \omterm{def:g9:landforms:dune}{وجه الانزلاق}. (ج) \omterm{def:g9:landforms:karst}{بالوعة}. (د) \omterm{def:g9:landforms:beach}{الانجراف الساحلي}.
(ه) \omterm{def:g9:landforms:karst}{الكارست}.
\end{solution}

\begin{solution}{exo:g9:landforms:3}
(أ) خطأ: القرنان يشيران مع اتجاه الريح؛ والأنف المستدير يواجه الريح.
(ب) صواب. (ج) خطأ: الريح تنقله حبيبةً حبيبةً، صعودًا على المنحدر اللطيف
ونزولًا على \omterm{def:g9:landforms:dune}{وجه الانزلاق}. (د) صواب.
\end{solution}

\begin{solution}{exo:g9:landforms:4}
(ب)، (د)، (ه)، (أ)، (ج).
\end{solution}

\begin{solution}{exo:g9:landforms:5}
(أ) من الغرب: فالقرون تشير مع اتجاه الريح. (ب) في الجانب الشرقي، بين
القرنين. (ج) $60 \div 6 = 10$ سنوات. (د) \omterm{def:g9:landforms:dune}{الكثيب} الأكبر يحوي \omterm{def:g3:stones-pebbles-sand:sand}{رملًا} أكثر،
ويجب أن يُحمل كله فوق الذروة كي يتقدم
\omterm{def:g9:landforms:dune}{الكثيب}.
\end{solution}

\begin{solution}{exo:g9:landforms:6}
(أ) $7 \times 1 = \qty{7}{m}$. (ب) $2 \div 12 = \frac{1}{6}$ من السنة.
(ج) سبع سنوات من التراجع في $\frac{1}{6}$ من السنة: $7 \times 6 = 42$
مرةً أسرع. (د) تحفر الأمواج أولًا حزًّا؛ ويبقى \omterm{def:g3:stones-pebbles-sand:rock}{الصخر} الذي فوقه بارزًا مدة
من الزمن، دون أن يتحرك؛ ثم يسقط البروز دفعة واحدة، فتقفز الحافة إلى
الوراء.
\end{solution}

\begin{solution}{exo:g9:landforms:7}
(أ) نحو الشرق. (ب) ثلاث ضربات صعودًا على \omterm{def:g9:landforms:beach}{الشاطئ} مائلة نحو الشمال
الشرقي، تتبع كلًّا منها ضربة عائدة مباشرة إلى أسفل نحو الجنوب: مسار
متعرج يتجه شرقًا. (ج) \omterm{def:g9:landforms:beach}{لسان رملي}، ينمو شرقًا عبر مدخل الخليج. (د) اللسان
يحمي الماء الذي خلفه من الأمواج؛ وفي الماء الهادئ، يستطيع \omterm{def:g3:stones-pebbles-sand:sand}{الوحل} الناعم
أن يترسب.
\end{solution}

\begin{solution}{exo:g9:landforms:8}
(أ) يأخذ ثاني أكسيد الكربون من الهواء ومن التربة. (ب) الحمض الضعيف يذيب
الحجر الجيري ببطء، لكنه لا يكاد يهاجم الغرانيت. (ج) لا يستطيع الماء أن
يدخل \omterm{def:g3:stones-pebbles-sand:rock}{الصخر} إلا عبر شقوقه وعلى طول طبقاته؛ فيوسع هذه الفتحات أولًا.
(د) يذوب \omterm{def:g3:stones-pebbles-sand:rock}{الصخر} الذي تحت السطح فتهبط الأرض ببطء؛ أو ينهار سقف
\omterm{def:g9:landforms:karst}{كهف}.
\end{solution}

\begin{solution}{exo:g9:landforms:9}
(أ) حين تلاقي القطرة هواء \omterm{def:g9:landforms:karst}{الكهف}، تطلق بعض ثاني أكسيد الكربون الذي فيها
فلا تعود قادرة على حمل كل الحجر الجيري المذاب، الذي يخرج منها كالسيتًا.
(ب) تنمو الصاعدة الكلسية من قطرات ترتطم بالأرض وتنتشر؛ فلا توجد قطرة
معلقة عند طرفها لتبني أنبوبًا.
(ج) $\qty{30}{cm} = \qty{300}{mm}$؛ $300 \div 0.1 = 3000 = 3 \times 10^3$
سنة.
\end{solution}

\begin{solution}{exo:g9:landforms:10}
(أ) $4600 \div 10^7 = 4.6 \times 10^{-4}$~m في السنة، أي \qty{0.46}{mm} في
السنة. (ب) $0.46 - 0.2 = \qty{0.26}{mm}$ في السنة. (ج) $\qty{1}{m} =
\qty{1000}{mm}$؛ $1000 \div 0.46 \approx 2200$، أي أسرع بنحو ألفي
مرة. (د) لأنه يستمر ملايين السنين: $\qty{0.26}{mm} \times
10^7 = 2.6 \times 10^6$~mm، أي نحو \qty{2.6}{km}.
\end{solution}

\begin{solution}{exo:g9:landforms:11}
(أ) الشكل: \omterm{def:g9:landforms:dune}{كثبان} هلالية وجه انزلاقها في الجانب الجنوبي الشرقي. العامل
والعملية: الريح، تهب من الشمال الغربي، \omterm{def:g5:water-shapes-the-land:deposition}{ترسيب}. المادة: \omterm{def:g3:stones-pebbles-sand:sand}{رمل} مفكك.
(ب) الشكل: \omterm{def:g9:landforms:cliff}{جرف} فوق رف مستوٍ ينكشف عند الجزر. العامل والعملية: \omterm{def:g2:water-on-earth:ocean}{البحر}،
\omterm{def:g5:water-shapes-the-land:erosion}{تعرية}؛ والرف مصطبة نحت أمواج تركها \omterm{def:g9:landforms:cliff}{الجرف} المتراجع. المادة: \omterm{def:g3:stones-pebbles-sand:rock}{صخر} أبيض، على
الأرجح طباشير، وهو \omterm{def:g3:stones-pebbles-sand:rock}{صخر} طري. (ج) الشكل: جدول يختفي في ثقب، \omterm{def:g6:water-underground:spring}{وينبوع} بعيد.
العامل والعملية: ماء يذيب \omterm{def:g3:stones-pebbles-sand:rock}{الصخر} ويجري تحت الأرض. المادة: حجر جيري؛
والهضبة
\omterm{def:g9:landforms:karst}{كارست}.
\end{solution}

\begin{solution}{exo:g9:landforms:12}
(أ) \omterm{def:g9:landforms:dune}{الكثيب} الصغير: $45 \div 6 = 7.5$ سنة؛ \omterm{def:g9:landforms:dune}{والكثيب} الأكبر: $45 \div 3 =
15$ سنة. (ب) $45 \div 100 = 0.45$ من السنة، أي نحو خمسة أو ستة أشهر.
(ج) مثلًا، زرع الأعشاب على \omterm{def:g9:landforms:dune}{الكثبان}: فجذورها تمسك \omterm{def:g3:stones-pebbles-sand:sand}{الرمل}، فلا تعود الريح
قادرة على أخذه.
\end{solution}

\begin{solution}{pb:g9:landforms:1}
\textbf{1.} $46.0 - 30.4 = \qty{15.6}{m}$.

\textbf{2.} $15.6 \div 20 = \qty{0.78}{m}$ في السنة.

\textbf{3.} بين 2012 و2014: $41.0 - 36.2 = \qty{4.8}{m}$ في
سنتين، أي $4.8 \div 2 = \qty{2.4}{m}$ في السنة.

\textbf{4.} $2.4 \div 0.78 \approx 3$: أسرع بنحو ثلاث مرات. على الأرجح
شتاء عاصف، سقطت خلاله بروزات كبيرة.

\textbf{5.} \qty{78}{cm} في السنة؛ $0.78 \times 100 = \qty{78}{m}$ في
القرن.

\textbf{6.} نعم: \qty{0.78}{m} في السنة أقل قليلًا من متر في السنة.

\textbf{7.} $17 \div 0.78 \approx 22$ سنة.

\textbf{8.} يتراجع \omterm{def:g9:landforms:cliff}{الجرف} قفزات: فقد تحوي فترة سنتين انهيارًا كبيرًا، أو
لا تحوي أي انهيار. وعلى مدى 20 سنة، تتوازن القفزات والسنوات
الهادئة.

\textbf{9.} أمام الحافة $30.4 - 10 = \qty{20.4}{m}$ لتقطعها: $20.4 \div 0.78
\approx 26$ سنة، أي في عام 2050.

\textbf{10.} $2 \times 0.78 = \qty{1.56}{m}$ في السنة؛ $30.4 \div 1.56
\approx 19.5$ سنة بعد 2024، أي في عام 2043.

\textbf{11.} قد تنهار حافة \omterm{def:g9:landforms:cliff}{الجرف} في أي لحظة؛ والوقوف عليها، أو قربها،
خطر.

\textbf{12.} $30.4 \div 0.78 \approx 39$ سنة بعد 2024: أي في عام
\textbf{2063}.
\end{solution}
